% Generated draft: requires mathematical and source review.
\documentclass[11pt]{article}
\usepackage[margin=25mm]{geometry}
\usepackage{fontspec}
\setmainfont{Latin Modern Roman}
\usepackage{amsmath,amssymb,mathtools,unicode-math}
\setmathfont{Latin Modern Math}
\usepackage{enumitem,needspace,xurl,hyperref}
\setlength{\parindent}{0pt}
\setlength{\parskip}{5pt}
\setlength{\emergencystretch}{3em}
\newcommand{\N}{\mathbb{N}}
\newcommand{\Z}{\mathbb{Z}}
\newcommand{\Q}{\mathbb{Q}}
\newcommand{\R}{\mathbb{R}}
\newcommand{\C}{\mathbb{C}}
\newcommand{\F}{\mathbb{F}}
\newcommand{\A}{\mathbb{A}}
\newcommand{\PP}{\mathbb{P}}
\newcommand{\E}{\mathbb{E}}
\newcommand{\eps}{\varepsilon}
\newcommand{\dd}{\,\mathrm{d}}
\begin{document}
\section*{UnsolvedMath ID: 800394}
\textbf{Source problem code:} catalogue entry
\subsection{Definitions and mathematical statement}
For a map \(f\), let \(f^0=\mathrm{id}\) and let \(f^n\) denote its \(n\)-fold iterate for \(n\geq 1\). The forward orbit of a point \(c\) is
\[
\mathcal O_f(c)=\{f^n(c):n\in\Z_{\geq 0}\}.
\]
A map is postcritically finite (PCF) if every critical point has finite forward orbit; equivalently, each critical point \(c\) satisfies \(f^m(c)=f^n(c)\) for some integers \(0\leq m<n\).

The source uses \(\mathcal M_d^{\mathrm{crit}}\) for a critically marked moduli space. A critical marking specifies critical points \(c_i\) of the represented map. Equalities of the form
\[
f^m(c_i)=f^n(c_j),\qquad m,n\in\Z_{\geq 0},
\]
express relations between marked critical orbits. These equalities explain the terminology, but the excerpt does not define the full class of relations intended by its conclusion.

\textbf{Supplied conjecture (Baker--DeMarco, approximately 2011; source Conjecture 17).}
Let \(X\subseteq\mathcal M_d^{\mathrm{crit}}\) be an algebraic subvariety, and put
\[
S_X=\{x\in X: x\text{ represents a PCF map}\}.
\]
The proposed implication is
\[
\overline{S_X}^{\,\mathrm{Zar}}=X
\quad\Longrightarrow\quad
X\text{ is cut out by ``critical orbit relations.''}
\]
This is asserted for every algebraic subvariety \(X\) in the source's intended moduli space.

\textbf{Unresolved qualifications.} The supplied excerpt does not specify the base field, the allowed values of \(d\), the underlying class of maps, or the precise marking and equivalence conventions defining \(\mathcal M_d^{\mathrm{crit}}\). It also does not define which critical orbit relations are allowed or whether ``cut out'' means a set-theoretic, scheme-theoretic, or componentwise characterization. Consequently, the displayed implication preserves the source formulation, but its conclusion cannot be made fully precise from this excerpt alone.
\subsection{Short English statement}
The supplied Dynamical Andr{\'e}--Oort conjecture predicts that an algebraic subvariety containing a Zariski dense set of postcritically finite maps is defined by relations among critical orbits. The intended moduli space and the exact meaning of those defining relations require clarification.
\subsection{Sources}
\begin{itemize}
\item Local UnsolvedMath registry: permanent ID \(800394\), source problem code \texttt{catalogue entry}, upstream ID \(800394\), upstream revision \texttt{d371ee917ac9647fbb06cd6fc529deec6f3dda82}. The supplied registry labels the problem \texttt{open}; this is a source label, not an independently verified status assessment.
\item Benedetto, DeMarco, Ingram, Jones et al., \emph{Current Trends and Open Problems in Arithmetic Dynamics} (2018), Conjecture 17, identified in the supplied record as the Dynamical Andr{\'e}--Oort Conjecture. Supplied source URL: \url{https://arxiv.org/abs/1806.04980}. The record reports access on 2026-10-16 and extraction by \texttt{arxiv\_latex}; neither access nor source contents have been independently verified here.
\item The statement attributes the conjecture to Baker--DeMarco, approximately 2011, and supplies the reference labels \texttt{arxiv1211.0255} and \texttt{MR2884382}. Full bibliographic details for these references are absent from the supplied record. The registry's proposer field lists Benedetto, DeMarco, Ingram, Jones et al., with proposed year 2018; this is preserved separately from the statement's attribution.
\end{itemize}
\end{document}
